\chapter{Life stages}

\section*{Definition and Overview}
Life stages represent the sequential developmental phases of human life, spanning infancy, childhood, adolescence, adulthood, and older age. Each stage is characterised by distinct physiological, psychological, and social milestones, as well as unique nutritional, immunisation, and preventive health needs. Understanding these stages is fundamental to clinical practice, enabling healthcare providers to deliver age-appropriate care, anticipate developmental challenges, and implement timely screening and intervention strategies.

\section*{Epidemiology}
The demographic landscape of human life stages is shifting globally, with a marked increase in the proportion of older adults due to declining birth rates and increased life expectancy. In Australia, the population aged 65 years and over constitutes approximately 16\% of the total population, a figure projected to rise. The distribution of disease burden varies dramatically across life stages: infectious diseases, developmental disorders, and accidents predominate in childhood and adolescence, whereas chronic non-communicable diseases, such as cardiovascular disease, type 2 diabetes, and dementia, represent the primary health burden in middle and late adulthood.

\section*{Aetiology and Pathophysiology}
Progression through life stages is governed by a complex interplay of genetic programming, hormonal signals, and environmental influences that dictate cellular growth, maturation, and eventual senescence.
\begin{itemize}
    \item \textbf{Infancy and childhood growth:} Driven by growth hormone and thyroid hormones, facilitating rapid cellular hyperplasia and hypertrophic growth, along with neural myelination and synaptic pruning.
    \item \textbf{Pubertal maturation:} Triggered by the activation of the hypothalamic-pituitary-gonadal (HPG) axis, leading to pulsatile secretion of gonadotropin-releasing hormone (GnRH), which stimulates the release of luteinising hormone (LH) and follicle-stimulating hormone (FSH) to promote sexual maturation.
    \item \textbf{Adult physiological maintenance:} A period of homeostatic stability where cellular repair mechanisms balance cellular damage, though lifestyle and environmental factors gradually influence systemic health.
    \item \textbf{Senescence and cellular ageing:} Characterised by progressive cellular damage, telomere shortening, mitochondrial dysfunction, and decreased regenerative capacity of stem cells, leading to organ system decline.
    \item \textbf{Hormonal decline:} Ageing is accompanied by programmed endocrine changes, such as the cessation of ovarian function in menopause and gradual testosterone decline in andropause.
\end{itemize}

\section*{Clinical Features}
Each life stage presents with specific physiological markers and clinical indicators of normal development or potential pathology.
\begin{itemize}
    \item \textbf{Infancy (0-2 years):} Rapid weight and height gain, acquisition of gross and fine motor skills (e.g. sitting, walking), language acquisition, and closure of cranial fontanelles.
    \item \textbf{Childhood (2-12 years):} Steady physical growth, cognitive development, motor coordination refinement, dental eruption transitions, and social integration.
    \item \textbf{Adolescence (12-18 years):} Development of secondary sexual characteristics, growth spurts, emotional volatility, identity formation, and increased risk-taking behaviours.
    \item \textbf{Adulthood (18-65 years):} Peak physical performance followed by gradual decline, career establishment, reproductive events, and early onset of chronic risk factors (e.g. hypertension, hypercholesterolaemia).
    \item \textbf{Older age (65+ years):} Frailty, cognitive slowing, sensory impairment (presbyopia, presbycusis), decreased bone density, muscle wasting (sarcopenia), and reduced physiological reserve.
\end{itemize}

\section*{Differential Diagnosis}
Clinical presentations must be interpreted within the context of the patient's life stage, as identical symptoms can have vastly different implications depending on age.
\begin{itemize}
    \item \textbf{Developmental delay vs. normal variation:} Distinguishing pathologically delayed milestones in infants from benign variations in developmental trajectories.
    \item \textbf{Pubertal delay vs. constitutional delay:} Differentiating hypogonadotropic hypogonadism from constitutional delay of growth and puberty.
    \item \textbf{Cognitive decline vs. normal ageing:} Differentiating mild cognitive impairment or dementia from benign senescent forgetfulness.
    \item \textbf{Fatigue across life stages:} Assessing fatigue as a potential indicator of depression or growth spurts in adolescents, compared to endocrine disorders or malignancy in older adults.
    \item \textbf{Joint pain causes:} Differentiating growing pains or juvenile idiopathic arthritis in children from mechanical osteoarthritis or inflammatory arthropathies in older adults.
\end{itemize}

\section*{When to Seek Medical Care}
Individuals or caregivers should seek medical attention if there are persistent deviations from typical developmental milestones, unexplained physical changes, or symptoms of systemic illness.

\textbf{Call 000 (or local emergency services) for:}
\begin{itemize}
    \item Sudden loss of consciousness, unresponsive states, or seizures
    \item Severe chest pain, respiratory distress, or cyanosis in any age group
    \item High fever in neonates under three months of age accompanied by lethargy
    \item Sudden, unexplained weakness, numbness, or difficulty speaking (signs of stroke)
\end{itemize}

\section*{Diagnosis and Investigations}
Diagnostic approaches are heavily tailored to the specific life stage to ensure safety and clinical relevance.
\begin{itemize}
    \item \textbf{Developmental screening tools:} Utilising structured questionnaires like the Ages and Stages Questionnaire (ASQ) to assess milestones in young children.
    \item \textbf{Paediatric growth charts:} Plotting height, weight, and head circumference against standardised centile charts to monitor growth velocity.
    \item \textbf{Endocrine assays:} Measuring gonadotropins, sex steroids, growth hormone, or thyroid hormones to evaluate developmental anomalies or transition phases.
    \item \textbf{Geriatric assessments:} Performing Comprehensive Geriatric Assessments (CGA), including cognitive tests (e.g. MMSE, GPCOG) and frailty scores.
    \item \textbf{Screening investigations:} Age-specific screenings such as newborn bloodspot screening, cervical screening in young adults, bowel screening in older adults, and bone densitometry (DEXA) scans.
\end{itemize}

\section*{Management}
Healthcare management must be individualised, transitioning from parental-supported care in childhood to autonomous self-management in adulthood, and supportive multidisciplinary care in older age.
\begin{itemize}
    \item \textbf{Paediatric immunisation:} Delivering age-specific vaccinations according to national schedules (e.g. Australian National Immunisation Program) and monitoring growth.
    \item \textbf{Adolescent transition care:} Managing mental health, addressing risk-taking behaviours, and providing sexual health education.
    \item \textbf{Adult preventive medicine:} Focusing on cardiovascular risk assessment, metabolic health management, and routine screening programs.
    \item \textbf{Geriatric care optimisation:} Minimising polypharmacy, preventing falls, managing cognitive decline, and establishing advance care directives.
    \item \textbf{Nutritional optimisation:} Adjusting dietary intake to meet life stage demands, such as folate in pregnancy, vitamin D and calcium in older age, and adequate protein for growing children.
\end{itemize}

\section*{Complications}
\begin{itemize}
    \item Failure to thrive or developmental stagnation in early childhood.
    \item Psychological issues, substance abuse, or self-harm arising from unaddressed adolescent distress.
    \item Uncontrolled metabolic risk factors in mid-adulthood leading to premature cardiovascular events.
    \item Sarcopenia, severe frailty, loss of independence, and institutionalisation in late adulthood.
    \item Progressive cognitive decline and dementia leading to complete loss of self-care capacity.
\end{itemize}

\section*{Prognosis and Outcomes}
The prognosis and long-term health outcomes of individuals are strongly influenced by interventions in early life stages. Healthy habits established in childhood and adulthood significantly reduce the incidence of chronic diseases in later stages. With modern medical advances, timely screening, and active lifestyle management, individuals can achieve a high quality of life and maintain functional independence well into their older years.

\section*{Prevention and Self-Care}
\begin{itemize}
    \item \textbf{Balanced nutrition:} Eat a diet rich in whole foods, adjusting caloric intake to match metabolic demands at different life stages.
    \item \textbf{Regular physical activity:} Maintain muscle mass and joint flexibility with age-appropriate cardiovascular and resistance exercises.
    \item \textbf{Routine screenings:} Participate in national screening programs (e.g. bowel, breast, cervical) and annual health assessments.
    \item \textbf{Cognitive stimulation:} Keep the mind active in older age through reading, puzzles, and social engagement to preserve cognitive function.
    \item \textbf{Vaccination compliance:} Maintain up-to-date immunisations, including annual influenza, pneumococcal, and shingles vaccines in older age.
\end{itemize}

\section*{Sources and Further Reading}
The following reputable organisations provide further information on life stages and human development:
\begin{itemize}
    \item Healthdirect Australia. \textit{Life stages and health resources.} https://www.healthdirect.gov.au/
    \item Better Health Channel. \textit{Healthy ageing and child development.} https://www.betterhealth.vic.gov.au/
    \item World Health Organization. \textit{Ageing and life course health.} https://www.who.int/
    \item National Institute on Aging. \textit{Health information for older adults and ageing research.} https://www.nia.nih.gov/
    \item Raising Children Network. \textit{Australian parenting website from newborns to teens.} https://raisingchildren.net.au/
\end{itemize}
